\documentclass[a4paper,11pt]{report}
\usepackage[catalan]{babel}
\usepackage[utf8]{inputenc}
\usepackage[T1]{fontenc}
\usepackage{geometry}
\usepackage{hyperref}
\usepackage{listings}
\usepackage{color}
\usepackage{longtable}
\usepackage{graphicx}
\usepackage{array}
\usepackage{pmboxdraw}

\geometry{margin=2.5cm}
\definecolor{dkgreen}{rgb}{0,0.6,0}
\definecolor{gray}{rgb}{0.5,0.5,0.5}
\definecolor{mauve}{rgb}{0.58,0,0.82}

\lstset{frame=tb,
  language=Python,
  aboveskip=3mm,
  belowskip=3mm,
  showstringspaces=false,
  columns=flexible,
  basicstyle={\small\ttfamily},
  numbers=none,
  numberstyle=\tiny\color{gray},
  keywordstyle=\color{blue},
  commentstyle=\color{dkgreen},
  stringstyle=\color{mauve},
  breaklines=true,
  breakatwhitespace=true,
  tabsize=3,
  inputencoding=utf8
}

\title{Resum Entorn i Documentació del Projecte}
\author{Josep Ferriol Font}
\date{\today}

\begin{document}
\maketitle
\tableofcontents

\chapter{Resum de l'Estructura del Projecte}

El projecte \textbf{TFG-truc} implementa el joc de cartes Truc amb una arquitectura modular preparada per a l'entrenament d'agents de Reinforcement Learning (RL) utilitzant la llibreria RLCard.

\section*{Visió General}
L'objectiu principal és proporcionar un entorn robust per simular partides de Truc, permetre interacció mitjançant interfícies gràfiques o de consola, i entrenar agents intel·ligents amb models avançats com DQN i NFSP.

\section*{Estructura de Directoris}

\begin{itemize}
    \item \textbf{\texttt{entorn/}}: Nucli de la lògica del joc i l'adaptador per a RLCard.
    \begin{itemize}
        \item \texttt{game.py}: Motor lògic del joc. Gestiona els estats de la partida, les regles, la progressió de rondes i mans, i les apostes (Truc i Envit).
        \item \texttt{env.py}: Adaptador de l'entorn per a RLCard. Tradueix l'estat intern del joc a observacions numèriques i exposa l'API estàndard de l'entorn (\texttt{reset}, \texttt{step}).
        \item \texttt{cartes\_accions.py}: Fitxer de constants compartides. Defineix la baralla espanyola de 36 cartes i l'\textbf{espai d'accions} complet (19 accions: 3 jugar carta, 7 de aposta/pas i 9 de senyals).
        \item \texttt{rols/}: Implementacions dels rols del joc (\texttt{dealer.py}, \texttt{judger.py}, \texttt{player.py}).
    \end{itemize}

    \item \textbf{\texttt{models/}}: Definició de les arquitectures de Xarxes Neuronals.
    \begin{itemize}
        \item \texttt{xarxa\_truc.py}: Defineix el \texttt{CosMultiInput}, que utilitza xarxes convolucionals i denses per processar el mapa de cartes i el context del joc en paral·lel. També inclou \texttt{ModelPreEntrenament}.
        \item \texttt{xarxa\_unificada.py}: Un wrapper end-to-end creat per unificar el Feature Extractor preentrenat amb l'algorisme RL de RLCard (DQN/NFSP). Pot carregar pesos prèviament apresos i congelar-los o re-entrenar-los (\textit{fine-tuning}).
        \item \texttt{loader.py} i \texttt{adapters/}: Utilitats per la càrrega de models desats.
    \end{itemize}

    \item \textbf{\texttt{controlador/}} i \textbf{\texttt{vista/}}: Implementació del patró MVC per separar la lògica gràfica i de flux del joc.
    \begin{itemize}
        \item \texttt{controlador/}: Conté els gestors interactius com \texttt{model\_interactiu.py} o \texttt{controlador.py} per enllaçar usuaris humans o bots amb la partida.
        \item \texttt{vista/}: Interfícies per a la interacció. Inclou \texttt{vista\_consola.py} i el directori \texttt{vista\_desktop/} per a la UI visual.
    \end{itemize}

    \item \textbf{\texttt{entrenament/}}: Lògica d'aprenentatge.
    \begin{itemize}
        \item \texttt{entrenamentEstatTruc/}: Conté \texttt{preentrenar\_cos.py}, que extreu datasets sintètics de partides simulatòries per entrenar el "Cos" de la xarxa mitjançant Supervised Learning.
        \item \texttt{entrenamentsUnificats/}: El nucli del Reinforcement Learning. L'script \texttt{entrenaments\_unificats.py} aglutina la càrrega del Cos, els agents DQN i NFSP i els enfronta amb tècniques com \textit{Reward Scheduling}, \textit{Learning Rate Decay} i un pool d'oponents històrics per l'aprenentatge òptim.
        \item \texttt{entrenamentRLCard/}: Versions prèvies i documentació teòrica associada com \texttt{entrenaments\_RLCard.md}.
    \end{itemize}

    \item \textbf{\texttt{notebooks/}}: Fitxers Jupyter interatius d'anàlisi de resultats.
    \begin{itemize}
        \item \texttt{Comparativa\_Experiments.ipynb}: Usat per extreure les gràfiques comparant els modes \textit{scratch}, \textit{frozen} i \textit{finetune} de tots dos models, juntament amb la lliga simulada final.
    \end{itemize}

    \item \textbf{\texttt{demo.py}}: Script de demostració interactiu per jugar per consola de manera ràpida.
    \item \textbf{\texttt{doc/}}: Documentació tècnica en Markdown i \LaTeX.
\end{itemize}

\chapter{Documentació de \texttt{env.py}}

\section{Visió General}

\texttt{TrucEnv} és l'entorn que adapta la lògica del joc (\texttt{TrucGame}) a la interfície estàndard de \textbf{RLCard}. La seva funció principal és \textbf{transformar l'estat del joc} (diccionaris llegibles per humans) en \textbf{vectors numèrics} (observacions) que un agent de Reinforcement Learning pot processar. S'hi inclouen modificacions recents per tal que un \textit{wrapper} pla l'adapti als nous models.

\hrulefill

\section{Arquitectura}

\begin{verbatim}
                    observació, acció
  +------------+  <------------------>  +------------+
  |  Agent RL  |                        |  TrucEnv   |
  +------------+                        |  (env.py)  |
                                        +-----+------+
                                              | estat, step
                                              v
                                        +------------+
                                        |  TrucGame  |
                                        |  (game.py) |
                                        +--+--+--+---+
                                           |  |  |
                              +------------+  |  +------------+
                              v               v               v
                        +----------+   +-----------+   +-----------+
                        |TrucDealer|   |TrucJudger |   |TrucPlayer |
                        +----------+   +-----------+   +-----------+
\end{verbatim}

\texttt{TrucEnv} hereta de \texttt{rlcard.envs.Env}, que proporciona el bucle estàndard \texttt{reset()} $\rightarrow$ \texttt{step()} $\rightarrow$ \texttt{get\_payoffs()}.

\hrulefill

\section{Constructor \texttt{\_\_init\_\_}}

\begin{lstlisting}
TrucEnv(config)
\end{lstlisting}

\subsection{Paràmetres de configuració (diccionari \texttt{config})}

\begin{longtable}{|l|l|p{8cm}|}
\hline
\textbf{Clau} & \textbf{Defecte} & \textbf{Descripció} \\
\hline
\texttt{num\_jugadors} & 2 & Nombre de jugadors \\
\texttt{cartes\_jugador} & 3 & Cartes per jugador \\
\texttt{puntuacio\_final} & 24 & Punts per guanyar \\
\texttt{senyes} & False & Activar fase de senyals \\
\texttt{player\_class} & None & Classe dels jugadors (e.g. \texttt{HumanPlayer}) o \textbf{diccionari} \texttt{\{id: Classe\}} per barrejar tipus \\
\texttt{verbose} & False & Activar missatges de debug del joc \\
\texttt{allow\_step\_back} & False & Permetre desfer passos (heretat d'RLCard) \\
\texttt{seed} & - & Seed per reproducibilitat \\
\hline
\end{longtable}

\hrulefill

\section{Variables Internes de l'Entorn}

\subsection{Mapeig de Cartes i Senyals}

\begin{longtable}{|l|l|p{8cm}|}
\hline
\textbf{Variable} & \textbf{Tipus} & \textbf{Descripció} \\
\hline
\texttt{cartes} & \texttt{list[str]} & Llista de totes les cartes del joc (36 cartes) \\
\texttt{carta\_map} & \texttt{dict[str, int]} & Mapa \texttt{carta $\rightarrow$ index} per codificar cartes a one-hot \\
\texttt{signal\_map} & \texttt{dict[str, int]} & Mapa \texttt{senya $\rightarrow$ index} per codificar senyals a one-hot \\
\hline
\end{longtable}

\subsection{Dimensions de l'Espai d'Estat}

\begin{longtable}{|l|l|p{6cm}|}
\hline
\textbf{Variable} & \textbf{Valor} (exemple 2J, 3C, -senyes) & \textbf{Descripció} \\
\hline
\texttt{num\_cartes} & 36 & Cartes úniques a la baralla \\
\texttt{espai\_joc\_cartes} & 37 & \texttt{num\_cartes + 1} (inclou slot ``buit'') \\
\texttt{espai\_hist\_cartes} & 6 & \texttt{num\_jugadors $\times$ cartes\_jugador} (slots historial) \\
\texttt{espai\_senya} & 10 & \texttt{len(ACTIONS\_SIGNAL) + 1} (9 senyals + buit) \\
\texttt{espai\_hist\_senyes} & 0 o 6 & 0 si \texttt{senyes=False}, sino \texttt{num\_jugadors $\times$ cartes\_jugador} \\
\texttt{espai\_info\_publica} & 10 & Puntuacions + apostes + situació \\
\textbf{\texttt{state\_size}} & \textbf{343} (sense senyes) & Mida total del vector d'observació clàssic \\
\hline
\end{longtable}

\subsection{Espais per a RLCard}

\begin{longtable}{|l|p{10cm}|}
\hline
\textbf{Variable} & \textbf{Descripció} \\
\hline
\texttt{state\_shape} & \texttt{[[state\_size]] * num\_jugadors} --- dimensions de l'observació per jugador \\
\texttt{action\_shape} & \texttt{[[len(ACTION\_LIST)]] * num\_jugadors} --- 19 accions possibles \\
\hline
\end{longtable}

\hrulefill

\section{Mètode Principal: \texttt{\_extract\_state(state)}}

Transforma el diccionari d'estat del joc en un vector numèric per ser utilitzat en les xarxes. Al nou model es retorna com un diccionari subdividit per ser processat pel \texttt{CosMultiInput}.

\subsection{Estructura de l'Estat Observat (\texttt{obs\_cartes} i \texttt{obs\_context})}

\begin{itemize}
    \item \textbf{\texttt{obs\_cartes}}: Formato \texttt{(6, 4, 9)}, que és utilitzat freqüentment per un Feature Extractor tipus Convolutional Neural Network (CNN). Aporta informació visual sobre la posició i el tipus de cartes que té el jugador i s'han llançat.
    \item \textbf{\texttt{obs\_context}}: Vector de 17 característiques escalars o senyals simples incloent puntuació, nivells de Truc, etc.
\end{itemize}

\subsection{Detall Clàssic de cada secció}

\subsubsection{1. Mà del Jugador (one-hot $\times$ \texttt{cartes\_jugador})}

Per cada slot de carta a la mà:
\begin{itemize}
    \item Si la carta existeix $\rightarrow$ marca \texttt{1} a la posició de la carta dins dels 36 valors
    \item Si el slot està buit (ja s'ha jugat) $\rightarrow$ marca \texttt{1} a la posició 37 (simbol ``buit'')
\end{itemize}

\subsubsection{2. Historial de Cartes (one-hot $\times$ \texttt{espai\_hist\_cartes})}

Codifica cada carta jugada durant la mà. Mateixa codificació que la mà: one-hot de 37 posicions. Permet a l'agent recordar quines cartes s'han jugat i per qui.

\subsubsection{3. Historial de Senyals (one-hot $\times$ \texttt{espai\_hist\_senyes})}

Només actiu si \texttt{senyes=True}. Per cada slot de senyal:
\begin{itemize}
    \item Si existeix una senya $\rightarrow$ marca \texttt{1} a la posició corresponent (0-8)
    \item Si el slot està buit $\rightarrow$ marca \texttt{1} a la posició 10 (simbol ``buit'')
\end{itemize}

\subsubsection{4. Informació Pública (10 valors escalars)}

\begin{longtable}{|l|l|p{8cm}|}
\hline
\textbf{Posició} & \textbf{Contingut} & \textbf{Descripció} \\
\hline
0 & \texttt{puntuacio[0]} & Punts de l'equip 0 \\
1 & \texttt{puntuacio[1]} & Punts de l'equip 1 \\
2 & \texttt{estat\_truc.level} & Nivell actual del Truc (1,3,6,9,24) \\
3 & \texttt{estat\_truc.owner + 1} & Qui ha cantat Truc (0=Ningú, 1=J0, 2=J1) \\
4 & \texttt{estat\_envit.level} & Nivell actual de l'Envit (0,2,4,6,...) \\
5 & \texttt{estat\_envit.owner + 1} & Qui ha cantat Envit (0=Ningú, 1=J0, 2=J1) \\
6 & \texttt{fase\_torn} & Fase actual (0/1/2) \\
7 & \texttt{ma} & Qui és mà \\
8 & \texttt{comptador\_ronda} & Nombre de rondes completades \\
9 & \texttt{id\_jugador} & ID del jugador que observa \\
\hline
\end{longtable}

\subsection{Retorn de \texttt{\_extract\_state}}

\begin{lstlisting}
{
    'obs': {
        'obs_cartes': np.ndarray,      # Tensorial de cartes
        'obs_context': np.ndarray      # Vector contextual 1D
    },
    'legal_actions': OrderedDict,     # Accions legals com a OrderedDict
    'raw_obs': state,                 # Estat original (diccionari llegible)
    'raw_legal_actions': ['passar', 'apostar_truc', ...],  # Noms de les accions
    'action_record': self.action_recorder  # Historial d'accions (d'RLCard)
}
\end{lstlisting}

\hrulefill

\section{Altres Mètodes}

\subsection{\texttt{get\_payoffs()} $\rightarrow$ Recompenses Finals}

Delega a \texttt{TrucGame.get\_payoffs()}, que retorna la diferència de puntuació:
\begin{lstlisting}
[score[0] - score[1], score[1] - score[0]]
\end{lstlisting}
L'agent que guanyi tindrà payoff positiu, el perdedor negatiu. A entrenaments s'ajusta aquesta recompensa amb el factor Beta per penalitzar fortament les derrotes al tram final.

\subsection{\texttt{\_decode\_action(action\_id)} $\rightarrow$ Descodificar Acció}

Converteix un índex d'acció numèric al seu nom string:
\begin{lstlisting}
ACTION_LIST[action_id]  # e.g. 0 -> 'play_card_0', 4 -> 'apostar_truc'
\end{lstlisting}

\subsection{\texttt{\_get\_legal\_actions()} $\rightarrow$ Accions Legals}

Aquest mètode de l'entorn delega directament a \texttt{TrucGame.get\_legal\_actions()}, que és on realment es calcula la llista d'accions permeses. Retorna una \textbf{llista d'índexs} (enters) que representen les accions vàlides.

\section{Espai d'Accions Complet}

Definit a \texttt{cartes\_accions.py}:

\begin{longtable}{|l|l|l|}
\hline
\textbf{Índex} & \textbf{Acció} & \textbf{Categoria} \\
\hline
0 & \texttt{play\_card\_0} & Jugar carta \\
1 & \texttt{play\_card\_1} & Jugar carta \\
2 & \texttt{play\_card\_2} & Jugar carta \\
3 & \texttt{apostar\_envit} & Aposta \\
4 & \texttt{apostar\_truc} & Aposta \\
5 & \texttt{vull\_envit} & Resposta \\
6 & \texttt{vull\_truc} & Resposta \\
7 & \texttt{fora\_envit} & Resposta \\
8 & \texttt{fora\_truc} & Resposta/Retirada \\
9 & \texttt{passar} & Control de flux \\
10--18 & \texttt{senya\_*} & Senyals (9 tipus) \\
\hline
\end{longtable}

\hrulefill

\section{Flux d'Ús Típic}

\begin{lstlisting}
# 1. Configurar
config = {'num_jugadors': 2, 'cartes_jugador': 3, 'senyes': False, ...}
env = TrucEnv(config)

# 2. Inicialitzar partida
state, player_id = env.reset()

# 3. Bucle del joc
while not env.is_over():
    action = agent.triar_accio(state)  # L'agent escull
    state, player_id = env.step(action)

# 4. Obtenir recompenses
payoffs = env.get_payoffs()  # [+diff, -diff]
\end{lstlisting}

\hrulefill

\section{Casuístiques Importants}

\subsection{Informació Asimètrica}
Cada jugador només veu les seves pròpies cartes a \texttt{ma\_jugador}. Les cartes dels rivals \textbf{no} s'inclouen al vector d'observació. L'historial de cartes jugades (\texttt{hist\_cartes}) sí que és públic.

\subsection{Codificació One-Hot}
S'utilitza one-hot en lloc de valors ordinals perquè la xarxa neuronal no assumeixi una relació lineal entre les cartes. Cada carta és un vector binari independent.

\subsection{Slot ``Buit''}
L'extra posició (\texttt{num\_cartes + 1} = 37) serveix per marcar slots que encara no tenen carta. Sense aquest símbol, el vector no podria distingir ``no s'ha jugat cap carta'' de ``s'ha jugat la carta 0''.


\chapter{Documentació de \texttt{game.py}}

\section{Visió General}

\texttt{TrucGame} és la classe central que implementa tota la lògica del joc del Truc. Gestiona l'estat complet de la partida: jugadors, cartes, apostes (Truc i Envit), rondes, i puntuació global. Actua com a \textbf{motor del joc}, independent de la interfície d'entorn (\texttt{env.py}).

\section{Constructor \texttt{\_\_init\_\_}}

\begin{lstlisting}
TrucGame(num_jugadors=2, cartes_jugador=3, senyes=False, puntuacio_final=24, player_class=TrucPlayer, verbose=False)
\end{lstlisting}

\subsection{Paràmetres de configuració}

\begin{longtable}{|l|l|p{6cm}|}
\hline
\textbf{Paràmetre} & \textbf{Defecte} & \textbf{Descripció} \\
\hline
\texttt{num\_jugadors} & 2 & Nombre de jugadors \\
\texttt{cartes\_jugador} & 3 & Cartes repartides per mà \\
\texttt{senyes} & False & Activa la fase de senyals \\
\texttt{puntuacio\_final} & 24 & Punts per guanyar la partida \\
\texttt{player\_class} & \texttt{TrucPlayer} & Classe o \textbf{diccionari} \texttt{\{id: Classe\}} dels jugadors (permet partides mixtes Human vs Bot) \\
\texttt{verbose} & False & Si és True, imprimeix missatges de debug per consola \\
\hline
\end{longtable}

\hrulefill

\section{Variables Internes del Joc}

\subsection{Jugadors i Rols}

\begin{longtable}{|l|l|p{6cm}|}
\hline
\textbf{Variable} & \textbf{Tipus} & \textbf{Descripció} \\
\hline
\texttt{players} & \texttt{list[TrucPlayer]} & Llista d'instàncies de jugadors \\
\texttt{dealer} & \texttt{TrucDealer} & Gestiona la baralla i el repartiment \\
\texttt{judger} & \texttt{TrucJudger} & Determina guanyadors de rondes, envits i mans \\
\texttt{np\_random} & \texttt{RandomState} & Generador aleatori per reproducibilitat \\
\hline
\end{longtable}

\subsection{Control del Torn}

\begin{longtable}{|l|l|p{6cm}|}
\hline
\textbf{Variable} & \textbf{Tipus} & \textbf{Descripció} \\
\hline
\texttt{ma} & \texttt{int} & ID del jugador que és \textbf{mà} (comença jugant). Rota cada mà nova. \\
\texttt{current\_player} & \texttt{int} & ID del jugador que ha d'actuar ara \\
\texttt{turn\_player} & \texttt{int} & Jugador amb el torn ``real'' (es restaura després d'una aposta) \\
\texttt{turn\_phase} & \texttt{int} & Fase actual del torn: 0=Senyals, 1=Joc i Apostes \\
\texttt{response\_state} & \texttt{ResponseState} & Estat de resposta pendent \\
\hline
\end{longtable}

\begin{longtable}{|l|p{10cm}|}
\hline
\textbf{Valor} & \textbf{Significat} \\
\hline
\texttt{NO\_PENDING} (0) & No hi ha cap aposta pendent de resposta \\
\texttt{TRUC\_PENDING} (1) & S'ha cantat Truc i s'espera resposta (Vull / Fora / Re-apostar) \\
\texttt{ENVIT\_PENDING} (2) & S'ha cantat Envit i s'espera resposta (Vull / Fora / Re-apostar) \\
\hline
\end{longtable}

\noindent\textbf{Important:}
Mentre hi ha una resposta pendent, el jugador contrari \textbf{només} pot respondre a l'aposta (vull/fora/re-apostar). No pot jugar cartes ni fer altres accions.

\hrulefill

\subsection{Puntuació i Historial}

\begin{longtable}{|l|l|p{6cm}|}
\hline
\textbf{Variable} & \textbf{Tipus} & \textbf{Descripció} \\
\hline
\texttt{score} & \texttt{list[int]} & Puntuació global \texttt{[equip\_0, equip\_1]} \\
\texttt{hist\_cartes} & \texttt{list[tuple]} & Historial de cartes jugades: \texttt{(player\_id, carta)} \\
\texttt{hist\_senyes} & \texttt{list[tuple]} & Historial de senyals: \texttt{(player\_id, senya)} \\
\texttt{round\_counter} & \texttt{int} & Nombre de rondes completades a la mà actual \\
\texttt{cartes\_ronda} & \texttt{list[tuple]} & Cartes jugades a la ronda en curs \\
\texttt{ronda\_winners} & \texttt{list[int]} & Guanyadors de cada ronda (-1 = empat) \\
\texttt{payoffs} & \texttt{list[int]} & Recompenses finals per a l'entorn RL \\
\hline
\end{longtable}

\subsection{Estat del Truc (Aposta de cartes)}

\begin{longtable}{|l|l|p{6cm}|}
\hline
\textbf{Variable} & \textbf{Tipus} & \textbf{Descripció} \\
\hline
\texttt{truc\_level} & \texttt{int} & Nivell actual del Truc (punts en joc per cartes) \\
\texttt{truc\_owner} & \texttt{int} & Qui ha cantat l'última aposta de Truc (-1 = ningú) \\
\texttt{previous\_truc\_level} & \texttt{int} & Nivell anterior (per si diuen ``Fora'', saben quants punts guanya el rival) \\
\hline
\end{longtable}

\textbf{Seqüència d'escalada del Truc:}
\begin{verbatim}
1 (per defecte) -> 3 (Truc) -> 6 (Retruc) -> 9 (Val Nou) -> 24 (Joc Fora)
\end{verbatim}

\subsection{Estat de l'Envit (Aposta d'envit)}

\begin{longtable}{|l|l|p{6cm}|}
\hline
\textbf{Variable} & \textbf{Tipus} & \textbf{Descripció} \\
\hline
\texttt{envit\_level} & \texttt{int} & Nivell actual de l'Envit \\
\texttt{envit\_owner} & \texttt{int} & Qui ha cantat l'últim Envit (-1 = ningú) \\
\texttt{envit\_accepted} & \texttt{bool} & Si l'envit ha estat acceptat \\
\texttt{previous\_envit\_level} & \texttt{int} & Nivell anterior (per calcular punts si diuen ``Fora'') \\
\hline
\end{longtable}

\textbf{Seqüència d'escalada de l'Envit:}
\begin{verbatim}
0 (no cantat) -> 2 (Envit) -> 4 (Un mes) -> 6 (Dos mes) -> Tots (Falta Envit)
\end{verbatim}

\noindent\textbf{Regla ``Tots'' (Falta Envit):}
Si cap equip supera els 12 punts, val 24 (guanya partida). Si algun equip supera els 12, val \texttt{24 - puntuacio\_del\_lider}.

\hrulefill

\section{Mètodes Principals}

\subsection{\texttt{init\_game()} $\rightarrow$ Inicialitzar Partida}

\begin{enumerate}
    \item Crea els jugadors amb \texttt{player\_class}
    \item Crea el dealer i el judger
    \item Barreja i reparteix cartes
    \item Inicialitza totes les variables d'estat
    \item Retorna \texttt{(estat, current\_player)}
\end{enumerate}

\subsection{\texttt{step(action)} $\rightarrow$ Avançar el Joc}

El mètode principal. Rep una acció (int o string) i actualitza l'estat.

\subsubsection{Flux de decisions:}

\begin{samepage}
\begin{verbatim}
step(action)
      |
      +--[Resposta pendent?]
      |        |
      |   ENVIT_PENDING --> vull_envit | fora_envit | apostar_envit
      |   TRUC_PENDING  --> vull_truc  | fora_truc  | apostar_truc
      |
      +--[NO_PENDING - Torn normal]
               |
               +-- Fase 0 (Senyals) --> passar  (avanca a Fase 1)
               |                    --> senya_* (registra, avanca a Fase 1)
               |
               +-- Fase 1 (Apostes  --> apostar_envit --> ENVIT_PENDING
                   i Joc)           --> apostar_truc  --> TRUC_PENDING
                                    --> fora_truc     --> Rival guanya, reset ma
                                    --> play_card_X
                                         |
                                    Fi ronda? -> Fi ma? -> Fi partida?
\end{verbatim}
\end{samepage}

\subsection{\texttt{\_get\_return\_state()} $\rightarrow$ Retorn Intel·ligent}

Comprova si l'única acció legal és \texttt{passar}. Si és així, \textbf{executa el \texttt{passar} automàticament} per evitar torns trivials (important per a l'entrenament d'agents RL).

\subsection{\texttt{get\_state(player\_id)} $\rightarrow$ Construir Estat}

Retorna un diccionari amb tota la informació visible pel jugador.

\subsection{\texttt{get\_legal\_actions()} $\rightarrow$ Accions Legals}

Determina quines accions pot fer el jugador actual segons la situació:

\begin{longtable}{|l|p{8cm}|}
\hline
\textbf{Situació} & \textbf{Accions disponibles} \\
\hline
\texttt{ENVIT\_PENDING} & \texttt{vull\_envit}, \texttt{fora\_envit}, (+ \texttt{apostar\_envit} si \texttt{level $\le$ 6}) \\
\texttt{TRUC\_PENDING} & \texttt{vull\_truc}, \texttt{fora\_truc}, (+ \texttt{apostar\_truc} si \texttt{level < 24}) \\
Fase 0 (Senyals) & \texttt{passar} + totes les \texttt{senya\_*} \\
Fase 1 (Apostes i Joc) & \texttt{play\_card\_0}, \texttt{play\_card\_1}, ... (segons cartes disponibles), \texttt{apostar\_envit}\textsuperscript{1}, \texttt{apostar\_truc}\textsuperscript{2}, \texttt{fora\_truc} \\
\hline
\end{longtable}

\textsuperscript{1} Només si \texttt{envit\_level == 0} i \texttt{round\_counter == 0} (primera ronda, sense envit previ)

\textsuperscript{2} Només si el jugador actual no és el propietari de l'última aposta de Truc

\subsection{\texttt{\_reset\_hand\_state()} $\rightarrow$ Reset per Nova Mà}

Quan acaba una mà (per guanyador o per "Fora"):
\begin{enumerate}
    \item Avança \texttt{ma} al següent jugador
    \item Barreja i reparteix de nou
    \item Reseteja totes les variables de Truc, Envit, rondes i historials
\end{enumerate}

\subsection{\texttt{is\_ma\_over()} i \texttt{is\_over()} $\rightarrow$ Fi de Mà / Partida}

Distingeix els dos nivells d'acabament:
\begin{itemize}
    \item \texttt{is\_ma\_over()}: retorna \texttt{True} si la \textbf{mà actual} ha acabat (hi ha guanyador o s'han esgotat les rondes).
    \item \texttt{is\_over()}: retorna \texttt{True} si la \textbf{partida} ha acabat (\texttt{max(score) >= puntuacio\_final}).
\end{itemize}

\hrulefill

\section{Casuístiques Especials}
\subsection{Determinar guanyador de la mà}

El \texttt{TrucJudger.guanyador\_ma()} segueix la següent lògica completament escalable per a partides de 3, 5 o N rondes:
\begin{enumerate}
    \item \textbf{Registre de victòries:} Per cada ronda completada, l'equip guanyador suma +1. Si la ronda ha resultat en empat, \textbf{tots dos equips} sumen +1 a la seva victòria.
    \item \textbf{Majoria absoluta:} Si un equip guanya la \textbf{majoria} de rondes ($\ge$ meitat de les rondes + 1) i l'altre no $\rightarrow$ guanya immediatament.
    \item Si s'han jugat totes les rondes o tots dos assoleixen la majoria alhora per culpa d'un empat:
    \begin{itemize}
        \item L'equip amb més victòries guanya.
        \item Si estan totalment igualats: guanya qui hagi guanyat la \textbf{primera ronda no empatada}.
        \item Si absolutament tot són empats: guanya el \textbf{mà}.
    \end{itemize}
\end{enumerate}

\subsection{Passar automàtic}

\texttt{\_get\_return\_state()} detecta quan l'única acció legal és \texttt{passar} i l'executa automàticament. Això evita que un agent RL hagi de "decidir" una acció trivial, millorant l'eficiència de l'entrenament.

\subsection{Fases del torn}

\begin{samepage}
\begin{verbatim}
  Fase 0: Senyals
      |  passar / senya_*
      v
  Fase 1: Joc i Apostes
      |  play_card_X / (o post-resposta a envit/truc)
      v
  Nova Ronda (torna a Fase 0 si senyes=True, sino Fase 1)
\end{verbatim}
\end{samepage}

\noindent\textbf{Nota:}
Si \texttt{senyes=False}, la fase 0 se salta completament i es comença directament a la fase 1.
\end{document}